\chapter{Exemplo Numérico do Sistema de Transmissão}
\label{apendice:exemplo-transmissao}

Para tornar concretos os cálculos descritos na \autoref{subsec:propulsao-direcao-veiculo}, este apêndice percorre passo a passo o que o software do Raspberry Pi calcula quando o operador acelera a 60\% na 1ª marcha, partindo do repouso. Os valores usados são os mesmos do código e das equações \autoref{eq:funcao-transferencia} a \autoref{eq:tau-efetivo}.

\section*{Parâmetros da 1ª marcha}

Da \autoref{tab:gear-params}, os parâmetros da 1ª marcha são:

\begin{itemize}
	\item $K$ (limitador) = 20\% --- o PWM máximo que a 1ª marcha permite
	\item ideal\_low = 0\%, ideal\_high = 15\% --- faixa onde o motor opera com eficiência
	\item $\tau_{base}$ = 2,0~s --- constante de tempo na zona IDEAL
\end{itemize}

\section*{Cálculo do target PWM}

O operador pressiona 60\% do acelerador. O software converte esse valor para o PWM alvo:

\begin{equation*}
	\text{target\_pwm} = \frac{\text{throttle}}{100} \times K = \frac{60}{100} \times 20 = 12{,}0\%
\end{equation*}

O valor 12\% cai dentro da zona IDEAL (0--15\%). Como a zona é classificada pelo PWM atual do motor (não pelo target), e o PWM vai de 0\% até 12\% neste exemplo, ele permanece na zona IDEAL durante todo o trajeto. A 1ª marcha é adequada para esse nível de acelerador.

\section*{Classificação das zonas de eficiência}

O software calcula a margem da zona SUBOPTIMAL como 25\% da largura da zona IDEAL, com mínimo de 2\%:

\begin{equation*}
	\text{margem} = \max(\text{ideal\_width} \times 0{,}25,\; 2{,}0) = \max((15 - 0) \times 0{,}25,\; 2{,}0) = \max(3{,}75,\; 2{,}0) = 3{,}75\%
\end{equation*}

Os limites das zonas ficam:

\begin{equation*}
	\text{sub\_low} = \max(0 - 3{,}75,\; 0) = 0\% \qquad \text{sub\_high} = \min(15 + 3{,}75,\; 20) = 18{,}75\%
\end{equation*}

Na 1ª marcha, a zona SUBOPTIMAL esquerda é inexistente porque ideal\_low já é 0\% (início da faixa operacional). Só existem SUBOPTIMAL à direita (15--18,75\%) e POOR (18,75--20\%) acima da zona IDEAL.

O mapa completo da 1ª marcha é:

\begin{table}[!h]
\captionsetup{width=16cm}
\Caption{\label{tab:zonas-1a-marcha} Zonas de eficiência da 1ª marcha}
\IBGEtab{}{%
	\begin{tabular}{cccc}
		\toprule
		Faixa de PWM & Zona & $M_{zona}$ & $\tau_{eff} = \tau_{base} \times M_{zona}$ \\
		\midrule \midrule
		0\% a 15\% & IDEAL & 1,0 & $2{,}0 \times 1 = 2{,}0$~s \\
		\midrule
		15\% a 18,75\% & SUBOPTIMAL & 10,0 & $2{,}0 \times 10 = 20{,}0$~s \\
		\midrule
		18,75\% a 20\% & POOR & 25,0 & $2{,}0 \times 25 = 50{,}0$~s \\
		\bottomrule
	\end{tabular}%
}{%
\Fonte{elaborado pelo autor.}%
\Nota{Na 1ª marcha, a zona IDEAL começa em 0\%, então não há zonas SUBOPTIMAL nem POOR abaixo da zona IDEAL. A faixa operacional cobre continuamente 0--20\%.}%
}
\end{table}

Na zona IDEAL, o motor responde com $\tau_{eff} = 2$~s. Fora dela, o multiplicador $M_{zona}$ torna a resposta 10 a 25 vezes mais lenta. O efeito prático é que o motor ``engasga'' quando a marcha não é a certa para a faixa de PWM, forçando o operador a trocar.

\section*{Aceleração na zona IDEAL (0\% a 12\%)}

Com PWM = 0\% e target = 12\%, a zona é IDEAL e $\tau_{eff} = 2{,}0$~s. O software aplica a equação diferencial de 1ª ordem via integração de Euler (conforme \autoref{eq:ode-motor}):

\begin{equation*}
	\Delta y = \frac{K - y}{\tau_{eff}} \times \Delta t = \frac{12{,}0 - \text{PWM}}{2{,}0} \times 0{,}01
\end{equation*}

A cada iteração, o ``erro'' ($12{,}0 - \text{PWM}$) diminui, e o passo fica menor. A \autoref{tab:iteracoes-euler} mostra algumas iterações.

\begin{table}[!h]
\captionsetup{width=16cm}
\Caption{\label{tab:iteracoes-euler} Iterações da integração de Euler na zona IDEAL da 1ª marcha}
\IBGEtab{}{%
	\begin{tabular}{ccccc}
		\toprule
		Tempo (s) & PWM (\%) & Erro (12,0 $-$ PWM) & $\tau_{eff}$ (s) & $\Delta y$ (\%) \\
		\midrule \midrule
		0,00 & 0,000 & 12,000 & 2,0 & 0,0600 \\
		\midrule
		0,50 & 2,717 & 9,283 & 2,0 & 0,0464 \\
		\midrule
		1,00 & 4,721 & 7,279 & 2,0 & 0,0364 \\
		\midrule
		2,00 & 7,386 & 4,614 & 2,0 & 0,0231 \\
		\midrule
		4,00 & 10,374 & 1,626 & 2,0 & 0,0081 \\
		\midrule
		6,00 & 11,403 & 0,597 & 2,0 & 0,0030 \\
		\bottomrule
	\end{tabular}%
}{%
\Fonte{elaborado pelo autor.}%
}
\end{table}

A solução analítica da \autoref{eq:resposta-degrau}, contada a partir de $t_0 = 0$~s quando PWM~=~0\%, é:

\begin{equation*}
	y(t) = 12{,}0 \times \left(1 - e^{-t/2{,}0}\right)
\end{equation*}

Os marcos temporais confirmam as propriedades do sistema de 1ª ordem \cite{franklin2013sistemas}:

\begin{table}[!h]
\captionsetup{width=16cm}
\Caption{\label{tab:marcos-temporais} Marcos temporais da resposta ao degrau na 1ª marcha}
\IBGEtab{}{%
	\begin{tabular}{cccc}
		\toprule
		Tempo & PWM (\%) & \% do caminho (0\% $\rightarrow$ 12\%) & Propriedade \\
		\midrule \midrule
		$1 \times \tau = 2{,}0$~s & 7,585 & 63,2\% & Definição de $\tau$ \\
		\midrule
		$2 \times \tau = 4{,}0$~s & 10,376 & 86,5\% & --- \\
		\midrule
		$3 \times \tau = 6{,}0$~s & 11,403 & 95,0\% & Critério de convergência \\
		\midrule
		$5 \times \tau = 10{,}0$~s & 11,919 & 99,3\% & Praticamente no valor final \\
		\bottomrule
	\end{tabular}%
}{%
\Fonte{elaborado pelo autor.}%
}
\end{table}

Tempo total para atingir 95\% do target: $3\tau$ = \textbf{6,0~s}.

\section*{O que aconteceria na zona POOR?}

Se o operador estivesse na 1ª marcha com PWM em 19\% (zona POOR, entre 18,75\% e 20\%), tentando forçar o motor além da zona IDEAL sem trocar para a 2ª marcha, o $\tau_{eff}$ seria:

\begin{equation*}
	\tau_{eff} = 2{,}0 \times 25{,}0 = 50{,}0~\text{s}
\end{equation*}

Com target em 20\% (acelerador a 100\%) e PWM atual em 19\%, o erro é apenas 1\%. Mesmo assim, o passo de Euler seria:

\begin{equation*}
	\Delta y = \frac{1{,}0}{50{,}0} \times 0{,}01 = 0{,}0002\%~\text{por iteração}
\end{equation*}

Para comparação, na zona IDEAL com o mesmo erro de 1\% o passo seria $1{,}0 / 2{,}0 \times 0{,}01 = 0{,}005\%$ por iteração, 25 vezes mais rápido. O motor quase não responde na zona POOR, sinalizando que o operador deveria ter trocado para a 2ª marcha (cuja zona IDEAL começa em 12\%, sobrepondo a 1ª marcha em 12--15\%).

\section*{Desaceleração: soltar o acelerador}

Com PWM = 10\% e target caindo para 0\%, o software usa $\tau$ dividido por 3 (conforme descrito na \autoref{subsec:propulsao-direcao-veiculo}):

\begin{equation*}
	\tau_{decel} = \frac{\tau_{eff}}{3} = \frac{2{,}0}{3} = 0{,}667~\text{s}
\end{equation*}

O passo de Euler:

\begin{equation*}
	\Delta y = \frac{10{,}0}{0{,}667} \times 0{,}01 = 0{,}15\%~\text{por iteração}
\end{equation*}

A desaceleração é 3 vezes mais rápida que a aceleração. O motor desacelera sensivelmente em cerca de 2~s.

\section*{Desaceleração com freio a 100\%}

Se além de soltar o acelerador o operador pisar no freio a 100\%:

\begin{equation*}
	\text{boost}_{freio} = 1 + 9 \times \frac{100}{100} = 10{,}0
\end{equation*}

\begin{equation*}
	\tau_{decel} = \frac{\tau_{eff}}{3 \times \text{boost}_{freio}} = \frac{2{,}0}{3 \times 10} = 0{,}0667~\text{s}
\end{equation*}

\begin{equation*}
	\Delta y = \frac{10{,}0}{0{,}0667} \times 0{,}01 = 1{,}5\%~\text{por iteração}
\end{equation*}

A desaceleração com freio máximo é \textbf{30 vezes mais rápida} que a aceleração normal. O motor para em aproximadamente 0,07~s.
